% =====================================================================
%  THÈME 3 — Arrondir & opérations — CORRIGÉ — Niveau DIFFICILE
% =====================================================================
\documentclass[11pt,a4paper]{article}
\usepackage[utf8]{inputenc}
\usepackage[T1]{fontenc}
\usepackage{lmodern}
\usepackage[french]{babel}
\usepackage{amsmath,amssymb,mathtools}
\usepackage{icomma}
\usepackage{siunitx}
\sisetup{output-decimal-marker={,},per-mode=symbol}
\usepackage[version=4]{mhchem}
\usepackage{xcolor}
\usepackage[most]{tcolorbox}
\usepackage{enumitem}
\usepackage{fancyhdr}
\usepackage[a4paper,margin=2.2cm,headheight=26pt,headsep=8pt]{geometry}

\definecolor{primary}{RGB}{150,98,162}
\definecolor{primarydark}{RGB}{92,52,110}
\newcommand{\cs}{chiffre significatif}
\newcommand{\CS}{chiffres significatifs}

\pagestyle{fancy}\fancyhf{}
\fancyhead[L]{\small\textcolor{primarydark}{\textbf{Fiche 1} — Chiffres significatifs}}
\fancyhead[R]{\small\textcolor{primarydark}{\textsc{Corrigé · Difficile · Thème 3}}}
\fancyfoot[C]{\small\thepage}
\renewcommand{\headrulewidth}{0.8pt}
\renewcommand{\headrule}{\hbox to\headwidth{\color{primary}\leaders\hrule height \headrulewidth\hfill}}

\newtcolorbox[auto counter]{cor}[1][]{enhanced,breakable,boxrule=0.8pt,arc=2pt,
  colback=primarydark!5,colframe=primarydark,left=3mm,right=3mm,top=2mm,bottom=2mm,
  fonttitle=\bfseries,coltitle=white,
  title={Correction — Exercice~\thetcbcounter\ifx\relax#1\relax\relax\else\ ---\ \textmd{\itshape #1}\fi}}

\setlength{\parindent}{0pt}\setlength{\parskip}{3pt}

\begin{document}

\begin{center}
{\Large\bfseries\color{primarydark} Corrigé — Niveau Difficile}\\[1pt]
{\color{primary} Thème 3 : calculs en chaîne \& cas mixtes}
\end{center}
\vspace{2mm}

\begin{cor}[Masse molaire puis nombre de moles]
\begin{enumerate}[label=\alph*),leftmargin=1.6em,itemsep=1pt]
  \item $M(\ce{CO2})=M(\ce{C})+2\,M(\ce{O})=\num{12,01}+2\times\num{16,00}=\num{12,01}+\num{32,00}$.
        C'est une \emph{addition} : la donnée à $2$ décimales impose $2$ décimales :
        $\boxed{M(\ce{CO2})=\qty{44,01}{\gram\per\mole}}$ ($4$~CS).
  \item $\displaystyle n=\frac{m}{M}=\frac{\qty{5,00}{\gram}}{\qty{44,01}{\gram\per\mole}}
        =\num{0,113610}\ldots\ \si{\mole}$ (valeur non arrondie).
  \item C'est une \emph{division} : min des \CS\ $=3$ (venant de $m=\num{5,00}$, car $M$ en a $4$).
        Donc $\boxed{n\approx\qty{0,114}{\mole}}$. La masse $m$ ($3$~CS) fixe la limite.
\end{enumerate}
\end{cor}

\begin{cor}[Volume d'un cube]
\begin{enumerate}[label=\alph*),leftmargin=1.6em,itemsep=1pt]
  \item $V=a^3=\num{3,20}^3=\qty{32,768}{\centi\metre\cubed}$ (brut).
  \item Multiplication : $a$ a $3$~CS $\Rightarrow V$ a $3$~CS. Chiffre supprimé $6\,(\ge5)$ :
        $\boxed{V\approx\qty{32,8}{\centi\metre\cubed}}$.
  \item « À $4$~CS » : $\qty{32,768}{}\to\qty{32,77}{\centi\metre\cubed}$ (supprimé $8\ge5$).
        \textbf{Différence :} « arrondir à $n$~CS » est un geste \emph{mécanique} qu'on subit quand
        l'énoncé l'impose ; « donner le bon nombre de CS » découle des \emph{données} via la règle de
        l'opération (ici $3$~CS). Les deux ne coïncident pas forcément.
\end{enumerate}
\end{cor}

\begin{cor}[La catastrophe de la soustraction]
\begin{enumerate}[label=\alph*),leftmargin=1.6em,itemsep=1pt]
  \item $m_{\text{dépôt}}=\num{47,52}-\num{47,48}=\num{0,04}$. Les deux pesées ont $2$ décimales
        $\Rightarrow$ résultat à $2$ décimales : $\boxed{m_{\text{dépôt}}=\qty{0,04}{\gram}}$.
  \item Chaque pesée avait $\mathbf{4}$~CS ; la différence n'en a plus que $\mathbf{1}$. La
        soustraction de deux nombres très voisins fait s'effondrer la précision.
  \item Si l'on arrondissait dès une étape intermédiaire, l'erreur d'arrondi se propagerait et
        s'amplifierait (surtout avant une soustraction voisine). En gardant des \emph{chiffres de
        garde} et en n'arrondissant qu'au résultat final, on limite cette erreur cumulée.
\end{enumerate}
\end{cor}

\begin{cor}[Masse volumique d'un liquide]
\begin{enumerate}[label=\alph*),leftmargin=1.6em,itemsep=1pt]
  \item Masse du liquide (\emph{soustraction}, $2$ décimales chacune) :
        $m=\num{80,36}-\num{50,00}=\qty{30,36}{\gram}$ ($4$~CS).
  \item $\displaystyle \rho=\frac{m}{V}=\frac{\qty{30,36}{\gram}}{\qty{30,00}{\milli\litre}}
        =\num{1,012}\ \si{\gram\per\milli\litre}$. \emph{Division} : min $=4$~CS
        $\Rightarrow \boxed{\rho=\qty{1,012}{\gram\per\milli\litre}}$.
  \item $\rho=\qty{1,012}{\gram\per\milli\litre}\times\num{1000}=\qty{1012}{\gram\per\litre}
        =\num{1,012e3}\ \si{\gram\per\litre}$.
  \item Étape a) : règle des \textbf{décimales} (soustraction). Étape b) : règle des \textbf{\CS}
        (division).
\end{enumerate}
\end{cor}

\end{document}
